\documentclass[10pt]{article}
\usepackage[margin=0.28in,landscape]{geometry}
\usepackage{booktabs}
\usepackage{graphicx}
\usepackage{caption}
\pagestyle{empty}

\begin{document}

\begin{table}[t]
\centering
\caption*{\huge Current Experimental Summary}
\large
\setlength{\tabcolsep}{4pt}
\renewcommand{\arraystretch}{1.35}
\begin{tabular*}{\textwidth}{@{\extracolsep{\fill}}llcrrrrrl@{}}
\toprule
Experiment & Setting & Ckpt. & MAE & RMSE & CRPS & 90\% Cov. & 90\% Width & Boundary Risk \\
\midrule
NASA rollout curriculum
& Single split
& best
& 0.0410
& --
& 0.0771
& 0.9971
& 0.8305
& 5.63\% sample out-of-range \\

CALCE mixed-protocol
& CS2\_34--37 $\rightarrow$ CS2\_38
& best
& 0.1108
& 0.1841
& 0.0922
& 0.9298
& 0.6186
& 0.47\% negative samples \\

CALCE 1C same-protocol
& CS2\_35/36 $\rightarrow$ CS2\_38
& best
& 0.0756
& 0.1459
& 0.0821
& 0.9853
& 0.7441
& 0.74\% lower-90 below zero \\

Severson/MIT in-distribution
& 98/21/21 cell split
& best
& 0.0042
& 0.0155
& 0.0040
& 0.9961
& 0.0360
& 0.0\% negative / out-of-range \\
\bottomrule
\end{tabular*}

\vspace{1.2em}
\normalsize
\raggedright
Notes: MAE/RMSE are computed from Monte Carlo sample-mean rollouts when posterior samples are used.
CRPS denotes empirical sample-based continuous ranked probability score.
Coverage and width refer to the 90\% prediction interval over valid rollout targets.
All entries are single-split exploratory results and should be treated as interim evidence rather than final statistical conclusions.
\end{table}

\end{document}
